\documentclass[10pt,a4paper]{article}

\usepackage[T1]{fontenc}
\usepackage[utf8]{inputenc}
\usepackage{lmodern}
\usepackage[margin=1.6cm]{geometry}
\usepackage{amssymb}
\usepackage{enumitem}
\usepackage[hidelinks]{hyperref}
\usepackage{tabularx}
\usepackage{xcolor}
\usepackage{fontawesome5}
\usepackage{pdfrender}

\pagestyle{empty}
\setlength{\parindent}{0pt}
\setlength{\parskip}{0.45em}
\setlist[itemize]{
  leftmargin=1.4em,
  itemsep=0.2em,
  topsep=0.25em,
  parsep=0pt
}

\definecolor{heading}{HTML}{1F4E79}
\newcommand{\cvsection}[1]{%
  \vspace{0.8em}%
  {\large\bfseries\color{heading}#1}\par
  \vspace{-0.35em}%
  {\color{heading}\rule{\textwidth}{0.6pt}}\par
  \vspace{0.1em}%
}
\newcommand{\contactsep}{\quad$\blacklozenge$\quad}
% White fill with dark blue outline
\newcommand{\cvicon}[1]{\textpdfrender{TextRenderingMode=FillStroke,LineWidth=0.5pt,FillColor=white,StrokeColor=heading}{#1}\ }
\newcommand{\entry}[4]{%
  \begin{tabularx}{\textwidth}{@{}Xr@{}}
    \textbf{#1} & \textbf{#2} \\
    \textit{#3} & \textit{#4}
  \end{tabularx}
}

\begin{document}

\begin{center}
  {\LARGE\bfseries Sara Su Ee Tan}\par
  \vspace{0.35em}
  \cvicon{\faEnvelope}\href{mailto:stan.suee@gmail.com}{stan.suee@gmail.com}
  \qquad
  \cvicon{\faPhone}+44 7401 053196
  \qquad
  \cvicon{\faHome}London, UK
  \qquad
  \cvicon{\faPassport}UK (ILR) holder\par
  \vspace{0.2em}
  \cvicon{\faLinkedin}\href{https://www.linkedin.com/in/sara-se-tan1/}{linkedin.com/in/sara-se-tan1}
  \qquad
  \cvicon{\faGithub}\href{https://Sara-Tan.github.io}{Sara-Tan.github.io}
\end{center}

\cvsection{SUMMARY}

PhD-trained biomedical scientist turned AI Engineer, combining seven years of life sciences research with LLM research and multi-agent system development.
At the University of Sheffield, I researched genetics, mechanobiology, and biophysics, leading projects funded by Marie Skłodowska-Curie and the Wellcome Trust (first-author papers in \textit{Nature Communications} and \textit{Development}).

At Huawei Technologies R\&D UK, I research inference-time LLM reasoning (first-author \textit{NeurIPS TTCL 2026 workshop}) and build agentic systems for research intelligence and 3D scene reconstruction.

Across my academic career, I built six tools to automate life sciences workflows, including the open-source \textit{QuantifyPolarity} (200+ downloads).

I have authored several peer-reviewed papers and reviewed for \textit{Development}, \textit{Cell Reports}, NeurIPS workshops, and ICLR.

\cvsection{SKILLS}

{\centering\textbf{\textit{Technical \& Analytical}}\par}

\textbf{LLMs \& Agents:} Multi-agent architectures, RAG, MCP, inference-time optimization \textit{(LangGraph, Hugging Face, vLLM, CrewAI)}

\textbf{Life Sciences:} Developmental biology, mechanobiology, genetics, molecular biology, quantitative image analysis, scientific software development

\textbf{Programming \& Data:} Python (PyTorch, pandas, Matplotlib), MATLAB, Streamlit, Git; statistics and data analysis

{\centering\textbf{\textit{Strategic \& Business}}\par}

\textbf{Research Project Management:} Scientific research proposals, grant acquisition, experimental design, and execution

\textbf{Business Strategy:} MBA Essentials (London School of Economics), Strategic Management (Copenhagen Business School)

{\centering\textbf{\textit{Communication}}\par}

\textbf{Communication:} Scientific writing (reports, academic papers, and grants); presentations; teaching

\textbf{Languages:} English (full proficiency), Mandarin (bilingual proficiency), German (basic proficiency)

\cvsection{WORK EXPERIENCE}

\entry{Huawei Technologies R\&D UK}{London, UK}{AI Technical Engineer}{2025--present}
\begin{itemize}
   \item Conducted research on inference-time LLM reasoning, developing \textit{test-time self-distillation}: a decoding method that steers the model towards better answers using contrastive prompts, with no retraining/reward models. Improved performance on maths, coding, and graduate-level science benchmarks across model scales (first-author paper, \textit{NeurIPS 2026 Workshop}).
  \item Building a multi-agent workflow that performs 3D scene reconstruction from natural-language user queries.
  \item Designed and deployed internal agentic tools to scale the team's technology intelligence: a literature-review agent that maps research landscapes and identifies gaps, and a news agent that curates and analyzes daily industry developments, reducing manual research effort and informing strategic decisions.
  \item Led AI technology foresight, tracking frontier AI research and turning it into actionable recommendations that informed product and research priorities.
  \item Assessed existing product lines to identify AI capability gaps and opportunities, scoping new initiatives with clear objectives.
\end{itemize}

\entry{University of Sheffield}{Sheffield, UK}{Postdoctoral Research Fellow}{2021--2025}
\begin{itemize}
  \item Led a Wellcome Trust-funded project investigating how mechanical forces influence organ development, combining genetics, microscopy, and computational image analysis; published as first author in \textit{Nature Communications} (2025).
  \item Drove the full project lifecycle, from ideation and planning to analysis, publication, and funding proposals, ahead of schedule.
  \item Maintained the open-source \textit{QuantifyPolarity} software, writing documentation and tutorials, supporting users, and releasing upgrades.
  \item Supervised research technicians and postgraduate students, contributing to talent development and project success.
  \item Initiated and led a ``Coding for Biologists'' workshop series, improving programming and data analysis capabilities among biologists.
\end{itemize}

\cvsection{SELECTED PUBLICATIONS}

\begin{enumerate}[leftmargin=1.6em, itemsep=0.25em, topsep=0.25em]
  \item \textbf{Tan, S. E.} et al. (2026). ``Beam Search as Test-Time Self-Distillation via Counterfactual Contexts.'' \textit{NeurIPS Workshop on Towards Test-Time Continual Learning Agents}.
  \item \textbf{Tan, S. E.}, \& Strutt, D. (2025). ``Tissue shear as a cue for aligning planar polarity in the developing \textit{Drosophila} wing.'' \textit{Nature Communications}.
  \item Brittle, A., Warrington, S. J., Strutt, H., Manning, E., \textbf{Tan, S. E.}, \& Strutt, D. (2022). ``Distinct mechanisms of planar polarization by the core and Fat-Dachsous planar polarity pathways in the \textit{Drosophila} wing.'' \textit{Cell Reports}.
  \item \textbf{Tan, S. E.} et al. (2021). ``QuantifyPolarity, a new tool-kit for measuring planar polarized protein distributions and cell properties in developing tissues.'' \textit{Development}.
\end{enumerate}

\cvsection{EDUCATION}

\entry{University of Sheffield}{Sheffield, UK}{Doctor of Philosophy in Biomedical Science}{2020}
\entry{RWTH Aachen University}{Aachen, Germany}{Master of Science in Biomedical Engineering}{2015}

\cvsection{AWARDS \& SCHOLARSHIPS}

\begin{itemize}
  \item Women Academic Researcher Award (2022--2023)
  \item Best Research Poster (2019)
  \item Marie Skłodowska-Curie Actions ITN Scholarship (2017)
  \item European Erasmus Mundus Scholarship (2015)
  \item Huawei Scholarship (2014)
\end{itemize}

\phantomsection\label{sec:projects}
\cvsection{SELECTED PROJECTS}

\begin{itemize}
  \item \textbf{Research-Bot (2025):} Agentic AI tool for automated literature reviews.
  \item \textbf{QuantifyPolarity (2019):} Open-source image analysis software, adopted and cited by research institutions worldwide.
  \item \textbf{Automated Cell Tracking (2018):} Machine-learning GUI for live-cell segmentation, tracking, and counting.
  \item \textbf{Cystic Fibrosis Quantification (2016):} GUI for automated fibrosis quantification in lung-cell microscopy images.
  \item \textbf{Focal Adhesion Tracking (2015):} Automated tracking of protein morphology and dynamics over time.
  \item \textbf{Pathological Brain Detection (2014):} Computer-aided diagnostic GUI for classifying pathological brain images.
\end{itemize}

\end{document}
